% Part: propositional-logic

\documentclass[../../include/open-logic-part]{subfiles}

\begin{document}

\olpart{pl}{વિધાનાત્મક તર્કશાસ્ત્ર}

\begin{editorial}
  આ ભાગમાં શાસ્ત્રીય વિધાનાત્મક તર્કશાસ્ત્રની સામગ્રી છે. પ્રથમ પ્રકરણ
  પ્રમાણમાં પ્રાથમિક છે અને તેમાં માત્ર વ્યાખ્યાઓ અને પરિણામોની યાદી છે;
  ઘણી સાબિતીઓ કરી નથી, પણ અભ્યાસપ્રશ્નો તરીકે છોડી છે. પ્રમાણપદ્ધતિઓ
  અને પૂર્ણતા પ્રમેયની સામગ્રી પ્રથમ-ક્રમ તર્કશાસ્ત્રના ભાગમાંથી,
  ``FOL'' ટૅગને false રાખીને, સામેલ કરી છે. તેથી વિધેયો, પદો અને
  પરિમાણકોને લગતું બધું રહી જાય છે અને !!{structure}s~$\Struct{M}$ની
  વાતને બદલે !!{valuation}s~$\pAssign{v}$ની વાત આવે છે.

  આ ભાગને આગળ વધારીને વધુ વિગત ઉમેરવાની અને સત્યતાફલનલક્ષી પૂર્ણતા
  જેવા બીજા વિષયો તથા પરિણામો સામેલ કરવાની યોજના છે.
\end{editorial}

\olimport[syntax-and-semantics]{syntax-and-semantics}

\tagfalse{FOL}

\olimport[../first-order-logic/proof-systems]{proof-systems}

\iftag{prfSC}{%
  \olimport[../first-order-logic/sequent-calculus]{sequent-calculus}
}{}

\iftag{prfND}{%
  \olimport[../first-order-logic/natural-deduction]{natural-deduction}
}{}

\iftag{prfTab}{%
  \olimport[../first-order-logic/tableaux]{tableaux}
}{}

\iftag{prfAX}{%
  \olimport[../first-order-logic/axiomatic-deduction]{axiomatic-deduction}
}{}

\olimport[../first-order-logic/completeness]{completeness}

\tagtrue{FOL}

\OLEndPartHook

\end{document}
